\documentclass[11pt,a4paper]{article}
\usepackage[utf8]{inputenc}
\usepackage[T1]{fontenc}
\usepackage[spanish,es-nodecimaldot]{babel}
\usepackage{amsmath,amssymb}
\usepackage[round]{natbib}
\usepackage{booktabs}
\usepackage{graphicx}
\usepackage[margin=2.5cm]{geometry}
\usepackage{hyperref}
\usepackage{xcolor}
\hypersetup{colorlinks=true,linkcolor=blue!50!black,citecolor=blue!50!black,urlcolor=blue!50!black}
\usepackage{caption}
\captionsetup{font=small,labelfont=bf}

\title{\bfseries La abstención perfecta es conocimiento nulo\\
\large Dos puntos fijos estables de un objetivo de abstención autorreferencial\\
en un modelo de lenguaje chico con memoria versionada}

\author{Maximiliano Speranza\\
\small Investigador independiente, Buenos Aires, Argentina\\
\small \href{https://orcid.org/0009-0005-0413-8554}{ORCID 0009-0005-0413-8554}\\
\small \texttt{mrsperanza@frba.utn.edu.ar}}

\date{Septiembre de 2026}

\begin{document}
\maketitle

\begin{abstract}
\noindent
Un detector de alucinaciones tendría que predecir los errores de su propio modelo, así que el blanco
natural de entrenamiento es autorreferencial. Cada ejemplo se etiqueta según si el modelo
\emph{se equivocaría en caso de contestar}. Entrenamos ese blanco en un modelo de lenguaje de
$863{,}730$ parámetros con un archivo versionado co-entrenado, dejando fijas la arquitectura, los datos
y el presupuesto, y encontramos que el objetivo tiene \textbf{dos puntos fijos estables}. Nueve
unidades se reparten $5$ y $4$ entre ellos, y el reparto lo gobierna \textbf{cuánta recuperación tenía
ya el modelo cuando empezó el entrenamiento de abstención}. Seis unidades que arrancaron de cero se
reparten $2$ y $4$, mientras que tres que heredaron $12{,}000$ pasos de pre-entrenamiento de
recuperación (RECUP $0{,}77$ a $0{,}80$) alcanzaron el punto útil $3$ de $3$.

El degenerado es el hallazgo. Cuatro de nueve unidades convergen a un estado que puntúa
$\texttt{nose} = 1{,}0000$ y $\texttt{invento} = 0{,}0000$, \emph{perfecto} en las dos métricas que la
literatura de abstención informa, mientras su exactitud global queda en $0{,}4065$, que es exactamente
el piso trivial que se obtiene abstiéndose en todo. No son detectores que fallan. Son
\textbf{detectores perfectamente calibrados de un generador que no sabe nada}, y la calibración es
honesta, porque el modelo no recupera y lo informa correctamente. Su cabeza binaria colapsó a una
constante con $\mathrm{AUC}$ de $0{,}522$ a $0{,}578$ sobre su propio blanco, y esa constante se
ubica donde la predice el prior. Para la unidad más joven, el $\mathrm{logit}$ de la tasa de error base
vale $1{,}513$ contra una mediana medida de $1{,}507$.

La bifurcación se decide temprano y es nítidamente legible. En el paso $2500$, sobre el lote de
evaluación de $512$ muestras que el propio bucle de entrenamiento ya computa, las cinco unidades
sobrevivientes habían emitido $2$, $4$, $4$, $8$ y $1$ respuestas, y las cuatro degeneradas habían
emitido \textbf{cero}. La regla ``cero contra no cero en el paso $2500$'' separa \textbf{40 de 40}
corridas con cabeza binaria en doce familias de este repositorio, y una unidad sobrevivió con una
\emph{sola} muestra de $512$. La fase muda inicial pertenece al blanco autorreferencial. Sobre el
control pareado que comparte semilla \emph{y} inicialización, tres unidades con un blanco fijo e
independiente del modelo leen entre $0{,}06$ y $0{,}20$ contra $0{,}98$ y $0{,}99$ de sus gemelas
autorreferenciales.

El punto degenerado es \textbf{absorbente} y no meramente lento. Medir la recuperación dentro de una
misma unidad con muestras pareadas a lo largo de $4000$ pasos más da $-0{,}0021$, o sea $0{,}4\sigma$,
y no se mueve. Informamos además una corrección a nuestra propia primera lectura, en la que una
tendencia entre unidades que sobrevive a un control de ruido ($r = 0{,}986$) igual no era
extrapolable, porque la curva es cóncava y un ajuste lineal escondía una saturación.

La consecuencia práctica es una afirmación sobre métricas y no sobre arquitectura. \textbf{Cualquier
criterio de éxito que lea tasa de abstención y tasa de fabricación sin leer exactitud global queda
maximizado por un modelo que no aprendió nada}, y un objetivo autorreferencial aporta una dinámica de
entrenamiento que alcanza ese máximo por sí sola, en una fracción sustancial de las corridas, sin
aviso de ninguna de las dos métricas.
\end{abstract}

\vspace{0.5em}
\noindent\textbf{Palabras clave.} abstención; detección de alucinaciones; predicción selectiva;
objetivos autorreferenciales; dinámica de entrenamiento; soluciones degeneradas; modelos de lenguaje
chicos; memoria persistente

\section{La métrica que gana un modelo callado}
\label{sec:intro}

La literatura sobre abstención informa dos números. Uno es la tasa con la que un modelo se niega
correctamente a contestar preguntas que no puede contestar, y el otro es la tasa con la que fabrica.
Los dos son deseables, los dos están acotados, y un sistema que mejora en ambos se toma normalmente
como un sistema que mejoró.

Un modelo que se abstiene en todas las preguntas puntúa perfecto en los dos.

No es una observación sutil, y como punto abstracto tampoco es nueva. La literatura de predicción
selectiva siempre emparejó riesgo con cobertura, precisamente porque el riesgo solo se minimiza de
forma trivial \citep{elyaniv2010selective, geifman2017selective}. Lo que informamos acá es que la
esquina degenerada no solo es alcanzable en principio sino que es un \textbf{atractor de un objetivo
de entrenamiento natural}, que una fracción sustancial de las corridas cae en ella, y que una vez
adentro el estado es estable si se sigue entrenando.

El objetivo en cuestión es el que un detector de alucinaciones realmente quiere. Un detector sirve si
predice los errores \emph{de este modelo}, así que el blanco no debería ser ``¿existe la respuesta?''
sino ``¿me voy a equivocar si contesto?''. Ese blanco está disponible gratis durante el entrenamiento.
Se toma el propio \texttt{argmax} del modelo, se compara contra la etiqueta y se usa la comparación
como blanco binario con el gradiente cortado. Es lo natural, y en este banco es lo que produjo las
unidades de las que trata este trabajo.

Su defecto es estructural antes que incidental. La etiqueta es función del propio estado del modelo,
así que se mueve mientras el modelo aprende, y el mapa tiene más de una solución consistente. En la
inicialización el modelo se equivoca en casi todo, con lo cual el blanco es la constante~$1$ y la
cabeza aprende a predecir ``me voy a equivocar'', que es \emph{cierto}. Escapar exige que el generador
empiece a acertar antes de que la cabeza termine de asentarse en la constante. Es una carrera entre dos
escalas de tiempo de aprendizaje, y la sección~\ref{sec:mudas} muestra que tanto la semilla como la
cantidad de recuperación que el modelo ya tiene cuando arranca la carrera deciden cuál gana.

Los dos desenlaces son autoconsistentes, y eso es lo que vuelve estable al degenerado.

\begin{center}
\begin{tabular}{ll}
\toprule
punto fijo & qué afirma el modelo, y si es cierto \\
\midrule
útil & ``recupero bien, y sé cuándo voy a fallar'', y es cierto \\
\textbf{degenerado} & ``no recupero casi nada, y sé que no recupero casi nada'', y \textbf{también es cierto} \\
\bottomrule
\end{tabular}
\end{center}

El modelo degenerado no está mal calibrado. Es honesto, inútil e invisible para las dos métricas que
el campo informa.

\paragraph{Contribuciones.}
\begin{enumerate}
\item Caracterizamos el punto fijo degenerado de un objetivo de abstención autorreferencial y mostramos
que lo alcanzan $4$ de las $6$ unidades entrenadas desde cero, y ninguna de las $3$ que heredaron
pre-entrenamiento de recuperación (sección~\ref{sec:mudas}).
\item Mostramos que las unidades degeneradas \emph{no} son detectores rotos. Sus cabezas colapsaron al
prior, y lo verificamos cuantitativamente contra la constante en forma cerrada, no solo de manera
cualitativa (sección~\ref{sec:prior}).
\item Mostramos que la bifurcación queda decidida para el paso $2500$ y es separable con una regla de
cero contra no cero sobre datos que el bucle de entrenamiento ya produce, validada en $40$ corridas
(sección~\ref{sec:bifurcacion}).
\item Mostramos que el punto degenerado es absorbente y no lento, usando la primera medición dentro de
una misma unidad que se hace en esta línea de trabajo (sección~\ref{sec:absorbente}).
\item Corregimos una primera lectura de nuestros propios datos, en la que una tendencia real entre
unidades se extrapoló con la forma funcional equivocada (sección~\ref{sec:correccion}).
\end{enumerate}

\section{Trabajo relacionado}
\label{sec:related}

\paragraph{Predicción selectiva.} El marco de riesgo y cobertura
\citep{elyaniv2010selective, geifman2017selective} existe porque la abstención se cambia contra la
utilidad, y SelectiveNet \citep{geifman2019selectivenet} vuelve explícito ese intercambio entrenando
una cabeza de selección contra una restricción de cobertura. Nuestras unidades degeneradas son la razón
de ser de esa restricción, observada como desenlace de entrenamiento y no como argumento.

\paragraph{Aprender a decir ``no sé''.} Agregar un token \texttt{[IDK]} al vocabulario
\citep{cohen2024idk} y agregar una cabeza binaria separada son las dos interfaces dominantes. Comparamos
cuatro de ellas con presupuesto igualado en un trabajo previo \citep{speranza2026abstention} y
encontramos que la cabeza es superior en este régimen. Esa comparación usaba un blanco \emph{fijo},
``¿existe la respuesta?'', y ninguna de sus unidades entró en el estado que se describe acá. El presente
trabajo es lo que pasa cuando el blanco se vuelve autorreferencial.

\paragraph{Blancos autorreferenciales y bootstrapeados.} Entrenar un predictor contra etiquetas
derivadas del modelo actual es estándar en autoentrenamiento y pseudo-etiquetado, y el modo de falla
por sesgo de confirmación está bien documentado \citep{arazo2020pseudo}. Nuestro aporte no es que los
blancos autorreferenciales puedan colapsar, cosa esperable, sino la forma específica del colapso en el
contexto de la abstención. El estado colapsado es \emph{correcto}, puntúa perfecto en las métricas
estándar y por lo tanto esas métricas no lo detectan como falla.

\paragraph{Soluciones degeneradas y colapso.} Los métodos de aprendizaje de representaciones que
predicen sus propias salidas colapsan a constantes si no se les pone un remedio explícito
\citep{grill2020byol, chen2021simsiam}. El mecanismo acá es de la misma familia, un blanco que el
modelo puede satisfacer moviendo el blanco, pero la solución colapsada no es una representación
constante sino una \emph{decisión} constante, y esa constancia es exactamente lo que la hace puntuar
bien.

\section{El banco}
\label{sec:bench}

El banco es el que describe \citet{speranza2026abstention}. Repetimos acá lo necesario para leer los
resultados y remitimos ahí para el resto.

\paragraph{Modelo.} Un modelo de lenguaje entrenado \emph{desde cero}, sin ningún componente
pre-entrenado. Son $863{,}730$ parámetros, $3{,}5$~MB y un vocabulario cerrado de $242$ tokens
legibles por humanos. Una recurrencia de regla delta más un \textbf{archivo persistente co-entrenado}
con un sello de orden aprendido en la clave, leído por softmax e inyectado en el bloque~0. El archivo
persiste entre sesiones mientras que el estado recurrente se reinicia entre ellas.

\paragraph{Tarea.} Los hechos se enuncian en una sesión, se corrigen en otra y se consultan en una
tercera. La respuesta es un \textbf{token único}, así que la métrica es exacta y determinista, sin juez
LLM y sin parser interpretativo. El $40\,\%$ de las preguntas no tiene respuesta
($p_{\mathrm{nose}} = 0{,}4$), y la fracción empírica en la muestra de evaluación es $0{,}4065$.

\paragraph{Interfaz.} Todas las unidades de este trabajo usan la cabeza binaria separada, que fue la
que seleccionó la comparación previa. Es un escalar $a$ proyectado desde el mismo estado final que
alimenta el softmax de vocabulario, con $+129$ parámetros. El modelo se abstiene cuando $a > 0$.
Informar el colapso sobre la interfaz \emph{ganadora} es deliberado, porque así no queda como artefacto
de un mecanismo de abstención débil.

\paragraph{Métricas.} \texttt{vigente} es la exactitud sobre el valor actual \emph{tal como se emite},
o sea después de la decisión de la cabeza. \texttt{nose} es la tasa de abstención en preguntas sin
respuesta. \texttt{falsa\_abst} es la tasa de abstención en preguntas que sí la tienen. \texttt{invento}
es la tasa de respuestas que nombran contenido ausente del archivo. Agregamos dos instrumentos.

\begin{itemize}
\item \textbf{Exactitud global}, definida como (respuestas correctas más abstenciones correctas) sobre
$n$. Su piso trivial en este banco es $\mathbf{0{,}4065}$, que se obtiene abstiéndose en todo.
\item \textbf{RECUP}, la exactitud de recuperación \emph{ignorando la cabeza}, o sea el \texttt{argmax}
sobre valores en las preguntas que sí tienen respuesta, sin aplicar la decisión de abstención.
\texttt{vigente} es un compuesto de recuperación y decisión, y vale $0{,}0000$ siempre que $a > 0$ en
todas partes, sin importar qué recuperó el modelo. RECUP separa las dos cosas.
\end{itemize}

\paragraph{Protocolo de medición.} RECUP se mide con $n = 8000$ y una semilla de datos fija ($54321$)
en todos los checkpoints, de modo que las comparaciones son \textbf{pareadas}. Las preguntas son
idénticas y lo único que cambia es el modelo. Esto importa. Medida sobre seis semillas de datos con
$n=2000$, RECUP tiene $\sigma = 0{,}0135$, así que una diferencia no pareada arrastra
$\sigma \approx 0{,}0191$, que es más grande que los efectos de interés. Informamos $p_{\mathrm{flip}}$,
la fracción de preguntas en las que dos checkpoints difieren, y de ahí derivamos el desvío estándar
pareado en vez de suponer un valor.

\paragraph{Protocolo.} Cada campaña tiene un preregistro congelado con su SHA antes de la
implementación, un archivo de desvíos y reporte por unidad. \textbf{Los resultados nunca se informan
como un promedio solo.} La bimodalidad entre semillas es una propiedad medida de este banco, y la
sección~\ref{sec:mudas} es una consecuencia directa de eso.

\paragraph{Uso de un modelo de lenguaje.} Todo el código de experimentación, las campañas de
entrenamiento y la generación de informes de este trabajo se hicieron con la asistencia de un modelo
de lenguaje grande. El diseño del estudio, los preregistros y sus hashes congelados, la elección de los
controles y la verificación de cada número informado son del autor. El modelo no es autor y no tiene
responsabilidad sobre el trabajo.

\section{El objetivo}
\label{sec:objetivo}

Escribimos $\ell$ para los logits de valor, $y$ para la etiqueta y $a$ para el escalar de la cabeza. El
blanco autorreferencial es
\[
b \;=\; \mathbb{1}\!\left[\operatorname{sg}\!\big(\arg\max\nolimits_{v \neq \texttt{NOSE}} \ell_v\big) \neq y\right],
\qquad
\mathcal{L} \;=\; \mathrm{BCE}(a, b) \;+\; \mathrm{CE}_{\text{valor}},
\]
con $\operatorname{sg}$ el corte de gradiente. En las preguntas sin respuesta, $b = 1$ por
construcción. La entropía cruzada de valor se computa solo donde existe respuesta y se normaliza por
esa fracción, de modo que la escala de la pérdida no dependa de $p_{\mathrm{nose}}$.

De ahí se siguen dos propiedades directas que conviene enunciar antes de cualquier medición, porque son
las que los resultados confirman.

\paragraph{El gradiente de valor nunca queda condicionado por la cabeza.}
$\mathrm{CE}_{\text{valor}}$ no depende de $a$. Le pase lo que le pase a la cabeza, la recuperación
sigue recibiendo gradiente. El colapso que informamos más abajo no es entonces un caso de ``el modelo
deja de aprender porque dejó de contestar'', una hipótesis que sostuvimos y refutamos
(sección~\ref{sec:mudas}).

\paragraph{Si la cabeza ignora su entrada, su óptimo es el prior.} Una constante $a$ que minimiza la
$\mathrm{BCE}$ contra $b$ vale $\mathrm{logit}\,\mathbb{E}[b]$, y $\mathbb{E}[b]$ se computa en forma
cerrada a partir de RECUP.
\begin{equation}
\mathbb{E}[b] \;=\; p_{\mathrm{nose}} + (1 - p_{\mathrm{nose}})\,(1 - \mathrm{RECUP}),
\qquad p_{\mathrm{nose}} = 0{,}4065 .
\label{eq:prior}
\end{equation}
Es una predicción sin ningún parámetro libre, y la sección~\ref{sec:prior} la pone a prueba.

\section{Resultados}

Nueve unidades, semillas $0$ a $8$, con banderas idénticas. Cinco llegan a un estado útil y cuatro no.
Informamos por unidad en todo el trabajo.

\paragraph{Un confundido en el punto de partida, encontrado tarde e informado completo.} Las nueve
unidades \emph{no} son homogéneas, y nos dimos cuenta después de escrito el análisis que sigue. El
arnés de entrenamiento siembra la fase de abstención desde un checkpoint base cuando existe uno para
esa semilla. Solo existían tres bases (\texttt{n3\_s0}, \texttt{n3\_s1}, \texttt{n3\_s2}, cada una con
$12{,}000$ pasos de entrenamiento de recuperación pura), así que las semillas $0$ a $2$ entraron a la
fase de abstención recuperando ya a RECUP $0{,}7734$, $0{,}7970$ y $0{,}7725$, mientras que las
semillas $3$ a $8$ arrancaron de cero. El preregistro de la campaña afirmaba que solo variaba la
semilla, y el chequeo que se hizo en su momento comparó los \emph{configs} guardados, que no registran
si hubo siembra.

\begin{center}
\begin{tabular}{lccc}
\toprule
grupo & semillas & RECUP inicial & desenlace \\
\midrule
pre-entrenadas & 0, 1, 2 & $\approx 0{,}78$ & \textbf{3 de 3 útiles} \\
desde cero & 3 a 8 & $0$ & 2 útiles, \textbf{4 degeneradas} \\
\bottomrule
\end{tabular}
\end{center}

Decimos con todas las letras qué cuesta esto y qué no. \textbf{Cuesta la tasa.} ``Cuatro de nueve
semillas'' no es una tasa sobre una población homogénea, y la afirmación honesta es $0$ de $3$ con
pre-entrenamiento y $4$ de $6$ sin él. \textbf{No cuesta el fenómeno}, que se mide por unidad y no
depende del agrupamiento. El punto fijo degenerado existe, cuatro unidades están en él, sus cabezas se
ubican en el prior y el estado es absorbente.

Y \emph{fortalece} el mecanismo de la sección~\ref{sec:objetivo} en vez de debilitarlo. Si el desenlace
es una carrera entre el aprendizaje de la recuperación y el asentamiento de la cabeza, entonces un
modelo que llega recuperando ya ganó la carrera antes de que empiece, que es exactamente el patrón
observado. Que dos de las seis unidades desde cero lleguen igual al punto útil muestra que el
pre-entrenamiento no es necesario, solo fuertemente favorable.

\textbf{Dónde muerde esto y dónde no, dicho en las dos secciones afectadas.} La
sección~\ref{sec:exclusiva} se apoya en un control pareado que comparte semilla \emph{y} base, así que
no queda tocada, mientras que la tabulación amplia de $31$ contra $9$ sí queda confundida y no nos
apoyamos en ella.

\subsection{Cuatro de nueve unidades terminan mudas}
\label{sec:mudas}

\begin{table}[h]
\centering
\caption{Las nueve unidades en su propio punto de operación. RECUP es la recuperación ignorando la
cabeza, y las últimas tres columnas describen el comportamiento de la cabeza. Las unidades degeneradas
van en negrita.}
\begin{tabular}{lrrrrrr}
\toprule
unidad & paso & RECUP & $\mathrm{AUC}(a)$ & $a > 0$ & mediana de $a$ & rango de $a$ \\
\midrule
b3\_s0 & 26000 & 0.9996 & 0.9997 & 0.407 & $-7.634$ & 29.17 \\
b3\_s1 & 26000 & 0.9992 & 0.9999 & 0.407 & $-7.563$ & 29.86 \\
b3\_s4 & 19000 & 0.7984 & 0.8084 & 0.413 & $-0.157$ & 14.17 \\
b3\_s2 & 26000 & 0.7946 & 0.8163 & 0.392 & $-0.233$ & 14.23 \\
b3\_s5 & 6000  & 0.7314 & 0.7550 & 0.529 & $+0.069$ & 5.66 \\
\midrule
\textbf{b3\_s6} & 26000 & \textbf{0.3964} & \textbf{0.5734} & \textbf{1.000} & $+1.006$ & \textbf{3.12} \\
\textbf{b3\_s3} & 26000 & \textbf{0.3820} & \textbf{0.5784} & \textbf{1.000} & $+1.045$ & \textbf{2.83} \\
\textbf{b3\_s7} & 13500 & \textbf{0.3432} & \textbf{0.5220} & \textbf{1.000} & $+1.251$ & \textbf{1.06} \\
\textbf{b3\_s8} & 8000  & \textbf{0.3040} & \textbf{0.5458} & \textbf{1.000} & $+1.507$ & \textbf{0.85} \\
\bottomrule
\end{tabular}
\label{tab:nueve}
\end{table}

\paragraph{Qué dicen las métricas estándar sobre las filas en negrita.} Las cuatro se abstienen en
todas las preguntas. Por lo tanto \texttt{nose} vale $1{,}0000$, o sea que rechazan correctamente toda
pregunta sin respuesta. Y \texttt{invento} vale $0{,}0000$, porque nunca nombran contenido ausente del
archivo, ya que nunca nombran nada. En los dos números que la literatura de abstención informa, estas
son las mejores unidades de la tabla, estrictamente mejores que \texttt{b3\_s0}, que contesta y por eso
ocasionalmente se equivoca.

Su exactitud global es $0{,}4065$, el piso trivial, con cuatro decimales, en las cuatro.

\paragraph{Una hipótesis que sostuvimos y refutamos.} \texttt{vigente} es compuesta, así que
$\texttt{vigente} = 0{,}0000$ es compatible con una recuperación \emph{perfecta} escondida detrás de una
cabeza trabada. Eso esperábamos. Es falso. RECUP va de $0{,}30$ a $0{,}40$ en las unidades degeneradas
contra $0{,}73$ a $1{,}00$ en el resto. La recuperación está genuinamente dañada.

Pero no es cero, y eso importa para la interpretación. Con $\approx 0{,}35$ contra una tasa de azar por
debajo de $0{,}01$, estos modelos aprendieron una parte sustancial de la tarea y se quedaron a mitad de
camino. El estado degenerado no es un modelo sin entrenar. Es un modelo parcialmente entrenado cuya
cabeza informa correctamente que ese entrenamiento parcial es insuficiente.

\subsection{La cabeza colapsó al prior, y la constante está donde la pone la teoría}
\label{sec:prior}

Las cuatro cabezas degeneradas no discriminan. Su $\mathrm{AUC}$ sobre su propio blanco va de $0{,}522$
a $0{,}578$, contra $0{,}9997$ de \texttt{b3\_s0}. Tampoco están saturadas en $+\infty$, ya que el rango
entero entre los percentiles $1$ y $99$ de sus logits abarca de $0{,}85$ a $3{,}24$, contra $29{,}17$.
Son constantes ubicadas apenas por encima del umbral de decisión.

La ecuación~\eqref{eq:prior} predice \emph{cuál} constante, sin ningún parámetro ajustado.

\begin{table}[h]
\centering
\caption{Prior en forma cerrada contra el logit mediano medido. No se ajusta ningún parámetro.}
\begin{tabular}{lrrrrr}
\toprule
unidad & RECUP & $\mathbb{E}[b]$ & $\mathrm{logit}$ predicho & mediana medida de $a$ & diferencia \\
\midrule
b3\_s8 & 0.3040 & 0.8196 & $1.513$ & $1.507$ & $\mathbf{0.006}$ \\
b3\_s7 & 0.3432 & 0.7963 & $1.363$ & $1.251$ & $0.112$ \\
b3\_s3 & 0.3820 & 0.7733 & $1.227$ & $1.045$ & $0.182$ \\
b3\_s6 & 0.3964 & 0.7647 & $1.179$ & $1.006$ & $0.173$ \\
\bottomrule
\end{tabular}
\label{tab:prior}
\end{table}

El ordenamiento se reproduce en las cuatro y la unidad más joven cae a $0{,}006$ de la predicción sin
parámetros. Las unidades más viejas quedan apenas por debajo de su prior, que es lo que haría una
cabeza que empieza a discriminar, y de manera consistente con eso el rango de $a$ crece de forma
monótona con el paso a lo largo de las cuatro ($0{,}85$, luego $1{,}06$, luego $2{,}83$, luego
$3{,}12$). \textbf{La cabeza no está rota. Es óptima para un generador que se equivoca el $80\,\%$ de
las veces.}

\subsection{La bifurcación se decide en el paso 2500}
\label{sec:bifurcacion}

Toda unidad con este blanco pasa por una fase muda inicial, cosa esperable, porque en la
inicialización el blanco es la constante~$1$. Lo que cambia es si esa fase termina.

El bucle de entrenamiento ya evalúa sobre $8$ lotes de $64$, o sea $512$ muestras. En el paso $2500$
tenemos lo siguiente.

\begin{table}[h]
\centering
\caption{Respuestas emitidas sobre $512$ en el paso $2500$, contra el desenlace final.}
\begin{tabular}{lrrl}
\toprule
unidad & \texttt{abstencion} @2500 & respuestas emitidas & desenlace \\
\midrule
b3\_s2 & 0.9844 & 8 & útil \\
b3\_s0 & 0.9922 & 4 & útil \\
b3\_s5 & 0.9922 & 4 & útil \\
b3\_s1 & 0.9961 & 2 & útil \\
b3\_s4 & 0.9980 & \textbf{1} & útil \\
\midrule
b3\_s3 & 1.0000 & \textbf{0} & degenerada \\
b3\_s6 & 1.0000 & \textbf{0} & degenerada \\
b3\_s7 & 1.0000 & \textbf{0} & degenerada \\
b3\_s8 & 1.0000 & \textbf{0} & degenerada \\
\bottomrule
\end{tabular}
\label{tab:bifurcacion}
\end{table}

La separación es cero contra no cero, y \texttt{b3\_s4} sobrevive con una \textbf{sola muestra de
$512$} y llega a $\texttt{vigente} = 0{,}68$. Romper el silencio una vez parece ser el evento que
importa, que es lo que predice la carrera de la sección~\ref{sec:objetivo}. Una muestra en la que el
\texttt{argmax} acierta es una muestra cuyo blanco vale $0$, y eso es justamente lo que una cabeza
constante no tiene gradiente para producir.

\paragraph{Validación sobre el repositorio.} Aplicada a \textbf{todas} las corridas de este repositorio
con cabeza binaria de abstención que llegaron a $\geq 6000$ pasos y tienen checkpoint en el paso
$2500$, o sea $40$ corridas, doce familias, dos reglas de decisión y los dos blancos, la regla separa
\textbf{40 de 40}. Son $36$ predichas útiles y útiles, y $4$ predichas degeneradas y degeneradas.

\paragraph{La prueba más dura disponible, corrida después de formular la regla.} La regla se derivó
post hoc. Después llevamos las dos unidades degeneradas \emph{más viejas} al horizonte completo con las
banderas originales, preregistrando la predicción de que ninguna saldría de la abstención total.
Ninguna salió. \texttt{b3\_s3} y \texttt{b3\_s6} llegaron al paso $26000$ con
$\texttt{abstencion} \geq 0{,}999$ en todos los hitos registrados. No es una validación prospectiva,
porque la prueba propiamente prospectiva habría sido un diseño de barrido que no corrimos, pero es la
prueba más fuerte que permitían las unidades existentes y la regla la sobrevivió.

\paragraph{Una advertencia sobre el margen.} Una muestra en $512$ es un margen del $0{,}2\,\%$. Como
criterio operativo la \emph{tasa} de abstención es demasiado frágil, y la forma robusta de la misma
cantidad es $\min_i a_i$ sobre un lote grande, que es continua e independiente del tamaño del lote.

\subsection{La fase muda pertenece al blanco autorreferencial}
\label{sec:exclusiva}

\textbf{Esta es la sección que más toca el confundido de la siembra, así que enunciamos primero el
límite.} Toda corrida con blanco fijo de este repositorio usa las semillas $0$ a $2$, y esas son
exactamente las semillas para las que existe una base pre-entrenada. La tabulación amplia que sigue
queda entonces confundida con la inicialización, y la informamos como descripción y no como evidencia.

\begin{center}
\begin{tabular}{lrl}
\toprule
blanco & corridas & \texttt{abstencion} en el paso 2500 \\
\midrule
fijo (``¿existe la respuesta?'') & 31 & $0{,}02$ a $0{,}32$ \\
\textbf{autorreferencial (``¿me voy a equivocar?'')} & 9 & $\mathbf{0{,}98}$ a $\mathbf{1{,}00}$ \\
\bottomrule
\end{tabular}
\end{center}

\textbf{La afirmación se apoya en cambio en el control pareado, al que el confundido no llega.} Las tres
unidades \texttt{p3} comparten semilla, arquitectura, presupuesto \emph{y} base pre-entrenada con
\texttt{b3\_s0}, \texttt{b3\_s1} y \texttt{b3\_s2}, y difieren solamente en el blanco.

\begin{center}
\begin{tabular}{lcc}
\toprule
semilla & \texttt{abstencion} @2500, blanco fijo & \texttt{abstencion} @2500, autorreferencial \\
\midrule
0 & 0.086 & \textbf{0.992} \\
1 & 0.201 & \textbf{0.996} \\
2 & 0.063 & \textbf{0.984} \\
\bottomrule
\end{tabular}
\end{center}

Tres de tres, con la variable confundida fijada por construcción. El blanco autorreferencial mete al
modelo en la fase muda \emph{incluso cuando llega recuperando} a RECUP $\approx 0{,}78$, y el blanco
fijo no lo hace nunca. Las unidades \texttt{p3} terminan con RECUP de $0{,}80$ a $0{,}97$ y ninguna cae
en la banda degenerada.

El mecanismo es el que nombra la sección~\ref{sec:objetivo}. Un blanco fijo es observable desde el
paso~$1$ y discriminativo de entrada, así que su cabeza nunca pasa por la fase constante en la que la
solución degenerada está disponible.

\subsection{El punto degenerado es absorbente y no lento}
\label{sec:absorbente}

Que el estado sea un punto fijo o un cuello de botella lento cambia por completo lo que se sigue de él,
y los dos se distinguen midiendo la recuperación \emph{dentro de una misma unidad}, cosa que esta línea
de trabajo nunca había hecho, porque los checkpoints intermedios se sobrescriben y la \texttt{vigente}
registrada vale $0{,}0000$ en todos los hitos sin importar la recuperación.

Archivamos cada checkpoint parcial de \texttt{b3\_s3} desde $22000$ hasta $26000$ y medimos RECUP en
cada uno con muestras pareadas ($n=8000$, semilla $54321$), bajo un preregistro congelado de antemano
que pedía $\geq +0{,}0100$ y monotonía en $\geq 3$ de $4$ intervalos.

\begin{center}
\begin{tabular}{lrrrr}
\toprule
paso & 22000 & 23000 & 24000 & 25000 \quad 26000 \\
\midrule
RECUP & 0.3841 & 0.3778 & 0.3844 & 0.3812 \quad 0.3820 \\
\bottomrule
\end{tabular}
\end{center}

El total es $\mathbf{-0{,}0021}$ a lo largo de $4000$ pasos. Con $p_{\mathrm{flip}} = 0{,}1526$ medido,
el desvío estándar pareado es $0{,}0057$, así que el total vale $0{,}4\sigma$ y cada intervalo queda
entre $0{,}1$ y $1{,}2\sigma$, o sea \textbf{ruido alrededor de una línea plana}. La monotonía dio $2$
de $4$. El criterio de riesgo del preregistro se disparó y se aplicó su cláusula de abandono, con lo
cual se canceló una campaña de seguimiento de $70000$ pasos que estaba planificada.

\paragraph{Qué no separa esto, declarado de antemano.} Los últimos $4000$ pasos corren con la tasa de
aprendizaje en el piso de su cronograma coseno. ``El atractor es absorbente'' y ``nada se mueve con la
tasa de aprendizaje en el piso'' son las dos compatibles con una curva plana, y la campaña cancelada era
la que las habría separado. Por eso enunciamos el resultado como \emph{absorbente bajo este presupuesto
y este cronograma}, y no como absorbente en general. Esta es la principal pregunta abierta que deja el
trabajo.

\paragraph{Un intervalo apunta al revés.} \texttt{b3\_s6}, al que le quedaban solo $1000$ pasos, da
$+0{,}0080$ ($1{,}6\sigma$). El preregistro ya había dictaminado que un intervalo suelto no constituye
tendencia, y no la constituye, pero es la única medición de este trabajo que empuja en dirección
contraria y la informamos como tal.

\subsection{Una corrección a nuestra propia primera lectura}
\label{sec:correccion}

Nuestro primer análisis ajustó una recta a través de las cuatro unidades degeneradas ordenadas por paso
y extrapoló que RECUP llegaría a $0{,}70$ hacia los $84000$ pasos, o sea que el estado era lento y no
absorbente. Antes de publicar esa lectura se corrieron dos chequeos, y el segundo la dio vuelta.

\paragraph{El chequeo que pasó.} La medición inicial usaba $n=2000$, cuyo $\sigma = 0{,}0135$ es
comparable a la dispersión entre unidades, así que sospechamos que la tendencia era ruido de medición.
Volver a medir las cuatro con $n=8000$ da $r = 0{,}9864$ con un rango de $12{,}5\sigma$. \textbf{La
tendencia es real y la sospecha estaba equivocada.}

\paragraph{El chequeo que falló.} La tendencia es real pero no es lineal. Las pendientes por intervalo
decrecen de forma monótona, con $+0{,}0071$, $+0{,}0048$ y $+0{,}0014$ por cada $1000$ pasos entre
unidades sucesivas, y $-0{,}0005$ dentro de \texttt{b3\_s3}, y un ajuste con saturación da una suma de
cuadrados residual $12{,}6$ veces menor que el lineal. La extrapolación era un artefacto de ajustar una
recta a una curva cóncava.

\paragraph{Qué \emph{no} concluimos del ajuste con saturación.} Tiene tres parámetros para cuatro
puntos, y su propia predicción en el paso $26000$ ($0{,}4535$) le erra al valor medido ($0{,}3820$) por
$0{,}07$. No es evidencia. La afirmación de la sección~\ref{sec:absorbente} se apoya en la medición
dentro de una misma unidad, que no tiene confundido entre semillas, y el ajuste con saturación solo
explica por qué la extrapolación anterior estaba mal.

\section{Lo que no afirmamos}
\label{sec:noclaim}

\paragraph{Esto no es evidencia de que el blanco autorreferencial sea inútil.} Cinco de nueve unidades
llegan a un estado útil, dos de ellas con $\mathrm{RECUP} \geq 0{,}999$ y $\mathrm{AUC} \geq 0{,}9999$.
La afirmación es que el objetivo admite una segunda solución estable y que una minoría sustancial de
semillas la encuentra, no que el objetivo fracase.

\paragraph{Las unidades degeneradas no están mal calibradas.} Sería cómodo llamarlas rotas. No lo
están. Sus cabezas son óptimas para sus generadores, con $0{,}006$ de diferencia contra la predicción
sin parámetros en el caso más limpio. La falla es que estar bien calibrado sobre no saber nada no vale
nada, que es una afirmación sobre el objetivo y las métricas, no sobre la calibración.

\paragraph{No afirmamos que el modelo sepa cuándo no sabe.} Todo acá es supervisado. La cabeza se
entrena contra etiquetas, y en el caso autorreferencial contra etiquetas derivadas de una comparación
supervisada. Las unidades degeneradas en particular son una advertencia contra esa lectura, porque
exhiben la firma externa de saber cuándo no saben mientras no saben absolutamente nada.

\paragraph{La regla del paso 2500 es post hoc.} Se formuló después de observar los desenlaces, se
validó sobre $40$ corridas existentes y se puso a prueba en dos unidades llevadas al horizonte completo
bajo preregistro. Eso es evidencia retrospectiva fuerte más una confirmación dura. No es una validación
prospectiva y no la informamos como tal.

\section{Limitaciones}
\label{sec:limits}

\begin{itemize}
\item \textbf{El punto de partida no estuvo controlado, y lo encontramos tarde.} Las semillas $0$ a $2$
heredaron $12{,}000$ pasos de pre-entrenamiento de recuperación y las semillas $3$ a $8$ no, así que el
reparto del titular no es una tasa sobre una población homogénea (sección~\ref{sec:mudas}). Los
fenómenos por unidad no quedan afectados y el control pareado de la sección~\ref{sec:exclusiva} está
limpio, pero cualquier campaña futura sobre esta pregunta tiene que fijar la inicialización en todas las
unidades, y el arnés tiene que registrar si hubo siembra, cosa que el config guardado no hace y por eso
el chequeo original no lo vio.
\item \textbf{Escala.} Son $863{,}730$ parámetros, un idioma sintético de $242$ tokens, un solo nivel de
dificultad y $p_{\mathrm{nose}}$ fijo en $0{,}4$. Nada de esto transfiere automáticamente. Solo
señalamos que el punto fijo degenerado es una propiedad del álgebra del objetivo antes que del tamaño
del modelo, lo que vuelve la pregunta por la transferencia digna de hacerse, no respondida.
\item \textbf{El presupuesto y el cronograma están confundidos con la afirmación de que es absorbente},
como se declara en la sección~\ref{sec:absorbente}. Es el hueco más grande del trabajo.
\item \textbf{Nueve unidades.} El reparto $5$ y $4$ es un conteo y no una estimación de tasa. Lo
informamos como observado y no le adosamos un intervalo de confianza a una proporción salida de nueve
semillas.
\item \textbf{Dos unidades están incompletas.} \texttt{b3\_s7} y \texttt{b3\_s8} se detuvieron en los
pasos $13500$ y $8000$, después de que el criterio bloqueante de la campaña se volviera aritméticamente
inalcanzable. Su condición degenerada en esos pasos está medida y no extrapolada, pero no corrieron
hasta el horizonte, y el desvío está declarado en el repositorio.
\item \textbf{Una sola regla de decisión.} Todas las unidades usan la cabeza binaria separada. Si las
interfaces de token o de slot admiten el mismo punto fijo bajo un blanco autorreferencial es algo que
no está probado.
\item \textbf{Búsqueda bibliográfica dirigida, no revisión sistemática.}
\end{itemize}

\section{Conclusión}

Un detector de alucinaciones quiere predecir los errores de su propio modelo, así que el blanco natural
es autorreferencial. Ese blanco tiene dos puntos fijos estables, y el degenerado no es una patología
para sacarse de encima por ingeniería sino una solución correcta. Es un modelo que no recupera casi nada
y que informa, con exactitud, que no recupera casi nada. Cuatro de las seis unidades que empezaron
desde cero lo encontraron, y ninguna de las tres que empezaron recuperando.

Lo que hace que esto sea más que una curiosidad es que la solución degenerada es \textbf{invisible para
las métricas que el campo informa}. Puntúa $\texttt{nose} = 1{,}0000$ e $\texttt{invento} = 0{,}0000$,
mejor que las unidades que funcionan, mientras su exactitud global queda en el piso trivial con cuatro
decimales. Cualquier criterio armado con tasa de abstención y tasa de fabricación queda maximizado por
ella. El remedio no es sutil, y consiste en informar la exactitud global contra su piso trivial y
exigirla junto a cualquier métrica de abstención, pero hay que aplicarlo, y en este proyecto se aplicó
recién después de que una campaña ya estuviera diseñada alrededor de criterios que un modelo mudo habría
ganado.

La bifurcación es legible temprano y barata de leer. Cero respuestas emitidas en el paso $2500$ separó
$40$ de $40$ corridas, y una unidad sobrevivió con una muestra de $512$. Si esa regla generaliza, cosa
que este trabajo no establece, entonces la pregunta relevante para un objetivo de abstención
autorreferencial no es cómo hacer que un modelo se abstenga bien. Abstenerse bien es trivial, y un
modelo puede lograrlo no sabiendo nada. La pregunta es cómo salir del silencio sin empezar a inventar, y
eso se decide en los primeros pocos miles de pasos.

\section*{Disponibilidad de datos y código}

Todos los preregistros con sus hashes SHA, los archivos de desvíos, los scripts de análisis, los
resultados crudos por semilla y los informes generados automáticamente están archivados en el
repositorio del proyecto, \url{https://github.com/SperanzaMax/delta-archivo}, cuyo historial de commits
lleva las marcas de tiempo que establecen el orden entre preregistro y resultado.

\section*{Declaraciones}

\textbf{Financiamiento.} Ninguno. \textbf{Conflictos de interés.} Ninguno. \textbf{Aprobación ética.}
No aplica, no hay sujetos humanos ni animales. \textbf{Disponibilidad de datos.} Ver más arriba.
\textbf{Contribuciones de autoría.} Autor único.

\bibliographystyle{plainnat}
\begin{thebibliography}{9}

\bibitem[Arazo et~al.(2020)]{arazo2020pseudo}
Arazo, E., Ortego, D., Albert, P., O'Connor, N.~E., and McGuinness, K. (2020).
Pseudo-labeling and confirmation bias in deep semi-supervised learning.
\emph{IJCNN}.

\bibitem[Chen and He(2021)]{chen2021simsiam}
Chen, X. and He, K. (2021).
Exploring simple Siamese representation learning.
\emph{CVPR}.

\bibitem[Cohen et~al.(2024)]{cohen2024idk}
Cohen, R., Dobler, K., Biran, E., and de~Melo, G. (2024).
I don't know: Explicit modeling of uncertainty with an [IDK] token.
\emph{NeurIPS}.

\bibitem[El-Yaniv and Wiener(2010)]{elyaniv2010selective}
El-Yaniv, R. and Wiener, Y. (2010).
On the foundations of noise-free selective classification.
\emph{JMLR}, 11:1605--1641.

\bibitem[Geifman and El-Yaniv(2017)]{geifman2017selective}
Geifman, Y. and El-Yaniv, R. (2017).
Selective classification for deep neural networks.
\emph{NeurIPS}.

\bibitem[Geifman and El-Yaniv(2019)]{geifman2019selectivenet}
Geifman, Y. and El-Yaniv, R. (2019).
SelectiveNet: A deep neural network with an integrated reject option.
\emph{ICML}.

\bibitem[Grill et~al.(2020)]{grill2020byol}
Grill, J.-B. et~al. (2020).
Bootstrap your own latent: A new approach to self-supervised learning.
\emph{NeurIPS}.

\bibitem[Speranza(2026)]{speranza2026abstention}
Speranza, M. (2026).
Where abstention lives: A pre-registered four-way comparison of abstention interfaces in a small
language model with versioned memory.
\emph{Research Square}, DOI 10.21203/rs.3.rs-10839567/v1.

\end{thebibliography}

\end{document}
